On a growing positive tail, the lip label is unbounded while its raw
abscissa remains small. The next estimate controls the logarithmic
derivatives of the actual exterior map in that regime. Keeping its
absolute first gain will later allow composition with the independent
negative section map.

\begin{lemma}[Complete section derivatives on a growing cubic tail]
\label{deg:lem-raw-growing-section}
Fix a sufficiently small terminal cut \(\delta>0\). Let
\[
 \begin{gathered}
 M(x)=x^3+\mu_2x^2+\mu_1x+\mu_0,\qquad B_0(x)=1+b+x,\\
 Q_0(x)=1-A(v+x)+C(v+x)^2,\qquad \zeta_0(x)=C(v+x).
 \end{gathered}
\]
with bounded \(A,C\). Hold all coefficients fixed as the starting
point \(a\) varies. On an enlarged interval \(I\subset(0,d_0]\),
with \(d_0<\delta/2\), suppose uniformly for \(a\in I\)
\[
 |b|+|v|\le C_0a,\qquad |\mu_j|\le C_0a^{3-j}\ (j=0,1,2),
 \qquad M(x)\ge c_0x^3\quad(a\le x\le\delta).
\]
Use the actual positive graph
\begin{equation}
 z_x=\frac{z\{xB_0(x)+\zeta_0(x)z\}}{D_0(x,z)},\qquad
 D_0=M+Q_0z,\qquad z(a)=\ell_0M(a),
 \label{deg:weighted-raw-positive-graph}
\end{equation}
and an increasing regular terminal adapter \(\chi\), with uniformly bounded
\(C^3\) norm and inverse on its enlarged image. For fixed positive
\(\sigma\), label the input by \(b_+=-a^2/(2\sigma^2)\), and write
\(\Gamma_+(b_+)=\chi(z_\delta)\). After choosing \(\ell_0,d_0\)
small, this continuation and every separate affine normalization
\(\phi_+\) satisfy
\begin{equation}
 |(a\partial_a)^j\log\phi_+'|\le Ca(1+|\log a|)^4,
 \qquad j=1,2.
 \label{deg:raw-growing-section-jets}
\end{equation}
The exact positive factor \(\sigma^2\) in the first derivative is
retained; no bound on its reciprocal is required.
\end{lemma}
\begin{proof}
Set \(E=a\partial_a\) and \(m(x)=M(x)/x^2\). On \(a\le x\le2a\),
the coefficient bounds and the polynomial margin give
\[
 c_0a\le m(x)\le C_1a,\qquad
 |b|+|v|+|\mu_2|+a+|\mu_1|/a+|\mu_0|/a^2\le C_1a.
\]
Thus Lemma~\ref{deg:lem-spectral-exterior} applies with \(t=1\)
and comparison mass \(\rho=a\), including its near-lip strip entry.
This mass is not differentiated;
all polynomial coefficients and \(\sigma\) remain fixed. The adapter
is evaluated on its stated enlargement of the generated terminal image.
In the actual raw-section applications,
\eqref{deg:spectral-exterior-window}, with \(\rho=a\), and the
original fixed terminal-image margin admit the whole family after
choosing \(d_0\) within its original cutoff reserve. The stated
ordinary \(C^3\) adapter and inverse bounds give a uniform bound for
its two logarithmic-height rows on the enlarged terminal image.

The common spectral estimate has
\[
 \lambda=\frac{4\mu_2}{1+b}=O(a),\qquad
 \Theta=a+|\mu_1|/a+|\mu_0|/a^2\le C_1a.
\]
Its terminal coefficient \(\mathcal T\) is bounded on the stated
buffer, and \(\mathcal T(a/\delta)^{2-\lambda}\le Ca\), after
the permitted \(d_0\) shrink. Consequently
\[
 |E\log\Gamma_+'+\lambda|+|E^2\log\Gamma_+'|
 \le Ca(1+|\log a|).
\]
The physical gain is given exactly by
\eqref{deg:scalar-polynomial-lip-slope},
\begin{equation}\label{deg:weighted-complete-positive-derivative}
 \Gamma_+'=\sigma^2\chi'(z_\delta)\frac{z_\delta}{D_0(\delta,z_\delta)}
                 B_0(a)J_0(a)e^{I(a)},\qquad J_0=1+O(a),
\end{equation}
where \(J_0=K\) and \(I\) are the exact factors in that formula.
This retains the factor \(\sigma^2\) without a reciprocal bound.
Since \(\lambda=O(a)\), both uncorrected logarithmic rows satisfy
the stronger first-logarithmic-power bound. Frozen affine calibration
leaves them unchanged, proving \eqref{deg:raw-growing-section-jets}.
\end{proof}

\begin{lemma}[Weighted cubic tail on the actual sections]
\label{deg:lem-weighted-cubic-tail}
Fix a compact set of coefficients \(0\le U\le U_0\),
\(|G|\le G_0\), and use the field
\eqref{deg:compact-cubic-field}--\eqref{deg:compact-cubic-polynomials}.
On its root-free positive fibres choose \(M\) sufficiently large,
independently of \(\tauc\), and retain
\[
 u\ge2M,\qquad a=\tauc u\le d_0,\qquad r_0\le r<c.
\]
Here \(c\) is the simple negative root and the two moving sections are
\(Y=\ell_0P\), with \(\eta=\tauc Y\).  The exterior maps are the
actual continuations to the fixed raw sections, followed by the regular
input adapter and the divided finite-side output adapter, respectively.
These continuations use the nested domains constructed by
\eqref{deg:compact-raw-corridor}.
There are constants, uniform on these domains, such that, with
\[
 E=u\partial_u,\quad
 \mathcal D_+=\frac{P(u)}{uB(u)}\partial_u,\quad
 \mathcal D_-=\frac{P(-r)}{rB(-r)}\partial_r,\quad
 \epsilon_u=\tauc(1+u),\quad L_u=1+|\log\epsilon_u|,
\]
the exact action matching has the form
\begin{equation}
 A_+(u)-A_-(r)=\tauc\{\log P(u)-\log P(-r)+R_+(u)-R_-(r)\},
 \label{deg:weighted-split-action}
\end{equation}
where
\begin{equation}
 |R_+|+|R_-|\le C,\qquad
 |E^jR_+|\le C\epsilon_u,\qquad
 |\mathcal D_-^jR_-|\le C\tauc\quad(1\le j\le3).
 \label{deg:weighted-action-jets}
\end{equation}
The relative variation in \eqref{deg:weighted-relative-variation} obeys
\begin{equation}
 |\varrho|+\frac{|\mathcal D_+\varrho|+|\mathcal D_-\varrho|}{u+r}
       \le C\epsilon_uL_u^4.
 \label{deg:compact-rho}
\end{equation}
If \(b_+=-u^2/2\),
\(b_-=-r^2/2\), and \(\phi_\pm\) are the separately affine-normalized
actual section embeddings, then
\begin{equation}
 |E\log\phi_+'|+|E^2\log\phi_+'|
       \le C\epsilon_uL_u^4,
 \qquad
 |\mathcal D_-^j\log\phi_-'|
       \le C\tauc(1+|\log\tauc|)^4\quad(j=1,2).
 \label{deg:weighted-actual-section-jets}
\end{equation}
Every prime on \(\phi_\pm\) is taken in its own label \(b_\pm\).
The full composition \(\Pi=\phi_-\circ S\circ\phi_+^{-1}\), with
\(S(-u^2/2)=-r^2/2\), therefore satisfies
\begin{equation}
 \phi_+'\frac{P(u)}{u+r}\frac{\Pi''}{\Pi'}
       =3-\frac{U}{ur}+E_u,
 \qquad |E_u|\le C\epsilon_uL_u^4.
 \label{deg:compact-curvature}
\end{equation}
The inward action estimates and the two exterior estimates are obtained
separately on their actual domains.
\end{lemma}

\begin{proof}
We first continue the two inward graphs to the original centre and
differentiate their exact height match. We then estimate the two
exterior embeddings separately. The negative embedding requires a
root-collar calculation because its first variation vanishes there.
Substitution in the completed-curvature identity gives the final estimate.

Choose a small fixed \(\delta_c>0\) such that \(P\) is a positive unit
on \([-2\delta_c,2\delta_c]\), and so that
\(2\delta_c<r_0\).  On the root-free positive fibre, put
\[
 \begin{gathered}
 q=\frac{A_+(u)-A_+(v)}\tauc,\quad
 W=\frac{Y}{P(v)},\quad D(v)=\frac{P(v)}{vB(v)},\\
 H(v)=\frac{P'(v)}{vB(v)},\qquad
 \varpi(v)=\frac{C\tauc^3(G+v)P(v)}{vB(v)}.
 \end{gathered}
\]
For \(\delta_c\le v\le u\), graph division gives the exact equations
\begin{equation}
 v_q=-\tauc D(v),\qquad
 W_q=W\left\{-\frac{1+\varpi W}{1+QW}+\tauc H\right\},
 \qquad W(0)=\ell_0.
 \label{deg:weighted-inward-equations}
\end{equation}
In particular, \(\tauc H=O(\epsilon_u)\),
\(Q-1=O(\epsilon_u)\), and
\(\varpi=O(\epsilon_u^3)\).  Fix \(\ell_0\), then \(d_0\),
small enough that the braces and the derivative of the right-hand side
with respect to \(W\), for \(0\le W\le\ell_0\), are both less than
\(-c_1<0\).  A first-exit argument gives
\begin{equation}
 0<W(q)\le\ell_0e^{-c_1q}.
 \label{deg:weighted-kernel-decay}
\end{equation}

For unbounded \(u\), Euler differentiation keeps the moving entrance
under control. At fixed action clock \(q\),
\begin{equation}
 Ev=\kappa(u)D(v),\qquad
 \kappa(u)=\frac{u^2B(u)}{P(u)}\asymp u^{-1}.
 \label{deg:weighted-transport}
\end{equation}
The first three Euler derivatives of \(\kappa\) are \(O(u^{-1})\).
On \([M,u]\), \(D(v)=O(v^2)\), and each further application of
\(D(v)\partial_v\) increases polynomial degree by at most one.
On \([\delta_c,M]\), this vector field and its first four derivatives
are uniformly bounded, including when the positive factor approaches a
double root: they contain \(P\), rather than \(P^{-1}\).
Differentiating \eqref{deg:weighted-transport} inductively therefore gives
\[
 |E^j\log v|\le C_j\quad(1\le j\le4),\qquad
 |E^j(\tauc H)|+|E^j(Q-1)|+|E^j\varpi|
       \le C_j\epsilon_u\quad(0\le j\le3).
\]
Because the transported field contains \(P\), these bounds remain
uniform when the root-free positive minimum tends to zero.
Apply Lemma~\ref{deg:lem-stable-ratio-ledger} with
\(D=E,\kappa=\tauc,\rho=\epsilon_u,\alpha=0\) on
\(\delta_c\le v\le u\).
The needed clock rows follow from \(\partial_q=-\tauc D(v)\partial_v\):
on the growing part its degree growth is controlled by
\(\tauc v\le\epsilon_u\), and on the compact part its coefficients
are bounded. The strict diagonal chosen above gives
\begin{equation}\label{deg:weighted-kernel-jets}
 |E^jW(q)|\le C\epsilon_u(1+q)^Ne^{-c_2q},\qquad1\le j\le3.
\end{equation}

The nonlinear logarithmic correction on this tail is exactly
\[
 V_+(u)=\int_0^{q_c(u)}
       \frac{(Q-\varpi)W}{1+QW}\,dq,
 \qquad q_c(u)=\frac{A_+(u)-A_+(\delta_c)}\tauc.
\]
The fixed-clock and terminal formulas of the same lemma give \(|V_+|\le C\) and
\(|E^jV_+|\le C\epsilon_u\), \(1\le j\le3\): indeed,
\(|E^jq_c|\le C_j/\epsilon_u\) and
\(q_c\ge[A_+(2M)-A_+(M)]/\tauc\ge a_M/\tauc\).
Every boundary term is therefore at most
\(C\tauc^{-N}e^{-c_2a_M/\tauc}\).

For the negative inward tail replace \((P,B,Q,\varpi)\) by
\((P(-v),B(-v),Q(-v),-C\tauc^3(G-v)P(-v)/(vB(-v)))\).
At fixed \(q\), \(\mathcal D_-\) transports to
\(P(-v)[vB(-v)]^{-1}\partial_v\);
its coefficient rows in the braces are \(O(\tauc)\) on
\([\delta_c,c]\), including the root factor. The same lemma with
\(\rho=\tauc\) gives \(|V_-|\le C\),
\(|\mathcal D_-^jV_-|\le C\tauc\). Since \(r\ge r_0>2\delta_c\),
its complete terminal rows are exponentially small.

Continue both tails to \(\xi=0\) using
Lemma~\ref{deg:lem-compact-height-continuation}, with radial coordinate
\(v=|\xi|\) oriented from \(\delta_c\) to zero and
\(\kappa_c=\tauc\). Here \(P\ge p_c>0\), the linear inward
propagator is at most one, and the nonlinear coefficient has ordinary
height derivatives through order three bounded by \(C/\tauc\).
For the incoming heights \(y_\pm=Y(\pm\delta_c)\), the exact memory
\eqref{deg:stable-ratio-polynomial-memory} gives
\(y_\pm\le C\tauc^{-N}e^{-c/\tauc}\); its three
logarithmic phase rows grow at most polynomially in \(\tauc^{-1}\).
The cited continuation and composition rows therefore make the remaining
nonlinear correction exponentially small, with those phase derivatives.
The exact cocycle \eqref{deg:exact-action-cocycle} gives the actual
heights on the common section:
\begin{align*}
 \log Y_0^+(u)&=\log(\ell_0P(u))-A_+(u)/\tauc+R_+(u),\\
 \log Y_0^-(r)&=\log(\ell_0P(-r))-A_-(r)/\tauc+R_-(r).
\end{align*}
with \eqref{deg:weighted-action-jets}.  Equality of these heights is
exactly the physical matching condition.  This proves
\eqref{deg:weighted-split-action}; it also supplies a separated extension
\(R_0(u,r)=R_+(u)-R_-(r)\), whose genuinely mixed derivatives vanish.

The separated height formulas also determine the relative first
variation and its weighted derivatives. Write
\(P_+=P(u)\), \(P_-=P(-r)\), \(B_+=B(u)\), \(B_-=B(-r)\),
and use radial derivatives in
\(H_+=P_+'/(uB_+)\), \(H_-=P_-'/(rB_-)\).
Differentiating \eqref{deg:weighted-split-action} gives
\begin{equation}
 r_u=\frac{uB_+P_-}{rB_-P_+}(1+\varrho),\qquad
 1+\varrho=
 \frac{1-\tauc H_+-\tauc\mathcal D_+R_+}
      {1-\tauc H_--\tauc\mathcal D_-R_-}.
 \label{deg:weighted-relative-variation}
\end{equation}
Since \(\mathcal D_+=(P_+/(u^2B_+))E\), the numerator differs
from one by \(O(\epsilon_u)\), with the same Euler bound.  The
denominator differs from one by \(O(\tauc)\), with two bounded action
variations of that order.  This proves \eqref{deg:compact-rho} (even
with \(L_u^4\) omitted).  Both denominators are positive after shrinking
\(d_0\) and then \(\tauc\).

For the positive exterior map apply
Lemma~\ref{deg:lem-raw-growing-section} with
\[
 \sigma=\tauc,\quad a=\tauc u,\quad
 (\mu_2,\mu_1,\mu_0)=(3\tauc G,-\tauc^2U,\tauc^3),
 \quad b=2\tauc G,\quad v=\tauc G.
\]
For \(u\ge2M\), the coefficient bounds hold and the leading cubic
term gives \(M(x)\ge c x^3\) on \([a,\delta]\), uniformly on the
fixed compact block. In the translated coordinate \(x=v-\gamma\),
the original positive raw port is \(x_T=\delta-\tauc G\).
The proof of the raw lemma is uniform for terminal cuts in a fixed
neighborhood of \(\delta\); \(x_T\) is coefficient-only and is held
fixed under \(E\). Thus it applies at this actual port. The corridor
supplies its enlarged terminal image.
Since \(a\asymp\epsilon_u\), \eqref{deg:raw-growing-section-jets}
is the positive assertion of \eqref{deg:weighted-actual-section-jets}.

Put \(\lambda_G=12\tauc G/(1+2\tauc G)\).
For the negative exterior map, first restrict to
\(r_0\le r\le c-x_0\), with \(x_0>0\) fixed and small.
Here \(P_-(r)=P(-r)\) has a fixed positive margin. Put \(v=-\xi\),
\(B_-(v)=B(-v)\), \(Q_-(v)=Q(-v)\), and use the increasing
height clock \(s=\log\eta\). The fast invariant
\(H=\eta-v^2/2\) satisfies the exact equation
\[
 H_s=\tauc\mathscr F(H,e^s),\qquad
 \mathscr F=
 \frac{v\eta(B_--Q_-)/\tauc
       -C\tauc(G-v)\eta^2-vP_-(v)}
      {vB_-(v)-C\tauc^2(G-v)\eta}.
\]
On a fixed box with \(v\ge r_0/2\), the denominator is a positive
unit, \((B_--Q_-)/\tauc\) is smooth and bounded, and
\(v=\sqrt{2(e^s-H)}\) gives
\( |\partial_H^i\partial_s^j\mathscr F(H,e^s)|\le C\)
for \(i+j\le3\). The initial rows are
\[
 s_l=\log\{\tauc\ell_0P_-(r)\},\qquad
 H_l=-r^2/2+\tauc\ell_0P_-(r).
\]
The ordinary \(r\)-derivatives of \(s_l-\log\tauc\) are bounded.
The moving initial recursion of
\eqref{deg:scalar-moving-bell} therefore gives the three initial
variations of \(-r^2/2\), with errors \(O(\tauc)\).
Choose a fixed small \(\eta_*>0\). First exit and the three
triangular equations on
\(s_l\le s\le\log\eta_*\), whose length is
\(O(1+|\log\tauc|)\), give
\[
 H(\log\eta_*;r)=-r^2/2+
          O_{C^3_r}\bigl(\tauc(1+|\log\tauc|)\bigr).
\]
This is an actual outward passage: initially
\(D=\tauc P_-+Q_-\eta>0\), and at a possible first zero
\(D_s=Q_-\eta>0\). On the fixed positive-height continuation
to a sufficiently large fixed cut \(v=L>c+1\), the graph equation
is an \(O_{C^3}(\tauc)\) perturbation of \(\eta_v=v\). Hence
\[
 \eta_L(r)=\tfrac12(L^2-r^2)+
          O_{C^3_r}\bigl(\tauc(1+|\log\tauc|)\bigr).
\]
The cut and its buffered chord lie in the original corridor
\eqref{deg:compact-raw-corridor}. From this cut to the raw port use
the negative tail with source quartet
\((a,p,\gamma,s)=(-\tauc,\tauc U,\tauc G,\tauc)\).
For a fixed \(0<\nu<1\), \eqref{eq:N-bound} gives
\[
 \int_L^{\delta/\tauc+G}\|\mathcal N\|_{C^3}\,dv
 \le C\tauc\int_L^\infty
 (v^{-2+\nu}+v^{-3+\nu}+v^{-4+\nu}+v^{-5+\nu})\,dv
 \le C\tauc.
\]
After extraction of its exact first slope, the tail and its inverse
are \(\mathrm{id}+O_{C^3}(\tauc)\). Its physical slope is
\(d_{\rm c}=\tauc^{2-\lambda_G}U_{\rm c}\), with
\(U_{\rm c}\) a positive unit. The original ordinary \(C^3\)
observer and divided closing add \(O_{C^3}(d_{\rm c})\) by
\eqref{eq:exterior-normalization}. Calibrating value and slope at a
frozen reference in this compact interval gives
\(\widetilde\phi_-=\mathrm{id}+O_{C^3}(\tauc(1+|\log\tauc|))\)
in \(b_-=-r^2/2\). Since \(\mathcal D_-\) has bounded coefficients
and derivatives here, the two logarithmic slope rows have the same
bound. They are unchanged by a frozen positive affine calibration,
so the original \(\phi_-\) satisfies
\eqref{deg:weighted-actual-section-jets} on this interval.
Near the root, the first variation vanishes and the preceding compact
normalization no longer gives a bounded inverse. We extract its endpoint
power before estimating logarithmic derivatives.

The root equation gives the exact factorization
\[
 P_-(r)=(c-r)L_-(r),\qquad
 L_-(r)=r^2+(U/c+c^{-2})r+c^{-1},\qquad
 L_-(c)=c^2+U+2/c>0.
\]
Thus \(c\) and \(L_-\), with their required finite coefficient
derivatives, are uniform units on a fixed negative collar,
independently of the positive bottleneck. Put \(x=c-r\), \(y=\eta\).
Dividing by the
positive normal rate \(rB(-r)\) gives
\[
 \dot x=-\beta x p(x)-k_1(x)y,\qquad
 \dot y=y\{1+k_2(x)y\},\qquad
 \beta=\frac{\tauc L_-(c)}{cB(-c)}=O(\tauc),
\]
Here the coefficients are exactly
\[
 \begin{split}
 p(x)&=\frac{cB(-c)L_-(c-x)}{L_-(c)(c-x)B(-c+x)},\\
 k_1(x)&=\frac{Q(-c+x)}{(c-x)B(-c+x)},\qquad
 k_2(x)=-\frac{C\tauc^2(G-c+x)}{(c-x)B(-c+x)}.
 \end{split}
\]
In particular \(p(0)=1\), \(p,k_1\) are positive units and
\(k_2=O_{C^7}(\tauc^2)\) on a fixed box.
Use the normal-time construction
\eqref{eq:saddle-normal-time}--\eqref{eq:exact-saddle-invariant} in the
proof of Lemma~\ref{lem:actual-dulac-matching}, with
\(\rho=\beta\), \(h=-\beta xp\), \(c=-k_1\), \(k=k_2\),
and selected coefficient parameter \(\tauc\) at fixed \(U,G,A,C\).
The explicit coefficients have a bounded joint order-seven budget:
\(h=O_{C_x^7}(\beta)\), \(k_1\) is a positive unit, and their
selected \(\tauc\)-rows are bounded. The normal coefficient and its
parameter rows retain their height factors.
The projection at phase order five uses axis-flow order six and the
selected parameter row. For a fixed normal decay rate \(0<\lambda<1\),
its integrands have the majorant
\(C(1+T)^Ne^{-(\lambda-C\beta)T}\), integrable after
\(C\beta<\lambda/2\). Thus the exact coordinates \((\Xi,Z)\)
and invariant \(I=\Xi Z^\beta\) have five phase derivatives and one
intrinsic \(\tauc\)-derivative, as required for the divisions below.
In the following coefficient-dependent units, a subscript
\({\rm c}\) records dependence on \(\tauc\). On the moving curve
\(y=\tauc\ell_0P_-\), these coordinates have the form
\[
 \Xi_m=xu_{\rm c}(x),\qquad Z_m=\tauc xv_{\rm c}(x),\qquad
 I(x)=\theta_{\rm c} x^{1+\beta}A_{\rm c}(x),\quad
 \theta_{\rm c}=\tauc^\beta,\quad A_{\rm c}=u_{\rm c} v_{\rm c}^\beta,
\]
with positive \(C^4\) units \(u_{\rm c},v_{\rm c},A_{\rm c}\).
To check the coefficient differences before taking the root limit, put
\(Y_{\rm c}(x)=\tauc x\ell_0L_-(c-x)\). Then
\[
 \partial_{\tauc}u_{\rm c}(x)=
 \frac{(\partial_{\tauc}\Xi)(x,Y_{\rm c}(x))}{x}
 +\Xi_y(x,Y_{\rm c}(x))\ell_0L_-(c-x),\qquad
 v_{\rm c}=\ell_0L_-(c-x)e^{-J(x,Y_{\rm c}(x))}.
\]
The first numerator vanishes at \(x=0\); its parameter-valued fifth
phase row gives \(C^4\) after Taylor division. In the second term the
moving factor \(x\) has already cancelled, so \(\Xi_y\) needs only
four phase rows. The formula for \(v_{\rm c}\) removes the factor
\(\tauc x\) before differentiation. Integrating these intrinsic
parameter rows gives \(u_{\rm c}-u_0,v_{\rm c}-v_0=O_{C^4}(\tauc)\),
and hence \(A_{\rm c}-A_0=O_{C^4}(\tauc)\). Also
\(\theta_{\rm c}=1+O(\tauc|\log\tauc|)\).

At a fixed outer cut beyond \(-c\), the stopped trajectory has
positive height. The invariant is a regular increasing coordinate
there, including \(I=0\), as in
\eqref{eq:exact-saddle-invariant}. Choose its buffered invariant chord
first, then the smaller lip collar so \(I(x)\) lies strictly inside it.
Choose the upper horizontal boundary with invariant values above that
chord throughout the local strip. This is possible at \(\tauc=0\),
where \(\pi=c-\sqrt{r^2-2\eta}\), and the strict margin persists
for small \(\tauc\). Thus no orbit of the lip collar exits through
that boundary.
The outward speed starts positive. At a possible first zero of
\(\mathscr D=\tauc P_-+Q_-\eta\), one has
\(\dot r=0\) and \(\dot{\mathscr D}=Q_-\dot\eta>0\), excluding
loss of that speed. Increasing height and the invariant chord then give
the actual first hit of the outer cut at positive height, uniformly for
\(0<x\le x_0\). Continue this fixed chord to the
negative raw section, before either the regular output adapter or the
finite closing map. After extracting the intrinsic exact first slope, write
\[
 T(x)=t_{\rm st}+d_{\rm c} F_{\rm c}(I(x)),\qquad
 d_{\rm c}=\tauc^{2-\lambda_G}U_{\rm c}>0,
\]
where \(U_{\rm c}\) is a positive unit and \(F_{\rm c},F_{\rm c}^{-1}\)
have uniform \(C^4\) bounds on a fixed invariant chord. Here the
finite positive-height passage has \(O(\tauc)\) coefficient changes.
On its subsequent tail, \eqref{eq:N-bound}, integrated with the
frozen negative quartet, bounds the nonlinear normalized kernel by
\(O(\tauc)\) through four phase orders. Thus the intrinsic,
separately normalized maps
satisfy \(F_{\rm c}=F_0+O_{C^3}(\tauc)\); no convergence rate for
the extracted slope \(d_{\rm c}\) is used.

Since \(b_-=-r^2/2\), one has \(db_-/dx=r=c-x>0\).
Consequently the exact first derivative is
\[
 T_{b_-}'=d_{\rm c}\theta_{\rm c} x^\beta
                 K_{\rm c}(x,x^\beta),\qquad
 K_{\rm c}(x,z)=
 \frac{F_{\rm c}'(\theta_{\rm c} xzA_{\rm c}(x))
       \{(1+\beta)A_{\rm c}(x)+xA_{\rm c}'(x)\}}{c-x}.
\]
The function \(K_{\rm c}\) is a positive \(C^3\) unit on a fixed
rectangle \(0\le x\le x_0\), \(0\le z\le1\), and
\(K_{\rm c}(0,z)\) is independent of \(z\).
At \(\tauc=0\) the outer graph is
\(\eta(\xi)=(\xi^2-r^2)/2\). Its intrinsic, separately
normalized section map is affine in \(b_-\). The regular factors
therefore cancel exactly at \(z=1\), giving a positive constant
\(k_0\), independent of \(x\), for which
\[
 \|K_{\rm c}(\cdot,1)-k_0\|_{C^2}
       \le C\tauc(1+|\log\tauc|).
\]
The affine limiting map is what makes the logarithmic derivative small;
boundedness of the regular factors alone would not suffice.

To differentiate along \(z=x^\beta\), use
\(E_x=x\partial_x+\beta z\partial_z\). For
\(0<x\le x_0<1\),
\[
 x|1-x^\beta|\le C\beta.
\]
The uniform \(C^3\) bounds for \(K_{\rm c}\), its constancy in
\(z\) at \(x=0\), and the preceding cancellation imply
\[
 |E_x\log K_{\rm c}|+|E_x^2\log K_{\rm c}|
       \le C\tauc(1+|\log\tauc|)
       \quad\hbox{on }z=x^\beta.
\]
Indeed the differences of the \(x\)-derivatives between \(z\)
and \(1\) have factors \(1-z\), accompanied by \(x\) or
\(x^2\); all remaining terms contain \(\beta\).
Since
\[
 \mathcal D_-=-\frac{L_-(c-x)}{(c-x)B(-c+x)}E_x,
 \qquad E_x\log x^\beta=\beta,
\]
we obtain the first two action-logarithmic bounds for
\(T_{b_-}'\), including their one-sided limits as \(x\downarrow0\).
The derivative itself may vanish like \(x^\beta\).

Finally append the original increasing ordinary \(C^3\) output adapter
\(\psi\) and the divided closing map. Put
\[
 \Psi_{\theta,m}=\mathfrak f_{\theta,m}\circ\psi,
 \qquad \mathfrak f_{\theta,m}(q)=f_\theta(mq)/m,
 \qquad \Gamma_-=\Psi_{\theta,m}\circ T.
\]
Admission to the original raw and closing windows precedes these
compositions. Apply \eqref{phys:eq-observer-log} and
\eqref{phys:eq-closing-jets} with \(D=\mathcal D_-\) and primes in
\(b_-\). They bound the two logarithmic-slope corrections using only
the ordinary \(C^3\) norms and positive first derivative of
\(\Psi_{\theta,m}\), uniformly for \(0\le m\le m_0\).
The invariant formula retains the small gain:
\(|DT|+|D^2T|\le Cd_{\rm c}\), and
\(d_{\rm c}\le C\tauc^2\), since \(\lambda_G=O(\tauc)\).
Both additional logarithmic rows are therefore \(O(d_{\rm c})\),
using only the ordinary \(C^3\) adapter and closing bounds of
Lemma~\ref{phys:lem-c3-adapters}. Thus the actual complete output
satisfies \eqref{deg:weighted-actual-section-jets}; affine normalization
fixes its logarithmic rows.

We can now combine the inward match with both exterior estimates.
Apply Lemma~\ref{neg:lem-completed-curvature} with
\(v=r\), \(P_+=P(u)\), \(P_-=P(-r)\), and \(W=u+r\).
In \eqref{deg:weighted-exact-curvature-chain}, put \(R_B=B_+/B_-\).
The first term differs from
\((P'(u)/u+P'(-r)/r)/(u+r)=3-U/(ur)\) by
\(O(\epsilon_u)\).  The second is \(O(\epsilon_u)\), since
\(R_B-1=\tauc(u+r)/B_-\) and
\(r_u=O(P_-/(u^2r))\).  The third is controlled by
\eqref{deg:compact-rho}, and the last by
\eqref{deg:weighted-actual-section-jets}, since
\(P_+/(u^2(u+r))\) is bounded.  For the negative section term use
\[
 \frac{\tauc(1+|\log\tauc|)^4}{u+r}
       \le C\epsilon_u(1+|\log\epsilon_u|)^4.
\]
To see this uniformly, write \(\epsilon_u=\tauc(1+u)\), use
\(1+|\log\tauc|\le1+|\log\epsilon_u|+\log(1+u)\), and note that
\((1+\log(1+u))^4/(1+u)^2\) is bounded.
This proves \eqref{deg:compact-curvature} in the actual affine
input coordinate.
\end{proof}
